%%%%%%%%%%%%%%%%%%%%%%%%%%%%%%%%%%%%%%%%%%%%%%%%
% E.Bigeard - emmanuel.bigeard@upsti.fr
% http://s2i.bigeard.me
% CC BY-NC-SA 2.0 FR - http://creativecommons.org/licenses/by-nc-sa/2.0/fr/
%%%%%%%%%%%%%%%%%%%%%%%%%%%%%%%%%%%%%%%%%%%%%%%%
\documentclass[11pt]{article}

%%%%%%%%%%%%%%%%%%%%%%%%%%%%%%%%%%%%%%%%%%%%%%%%
% Package UPSTI_Document
%%%%%%%%%%%%%%%%%%%%%%%%%%%%%%%%%%%%%%%%%%%%%%%% 
\usepackage{UPSTI_Document}

%---------------------------------%
% Paramètres du package
%---------------------------------%

% Version du document (pour la compilation)
% 1: Document prof
% 2: Document élève
% 3: Document à publier
\newcommand{\UPSTIidVersionDocument}{3}

% Affichage personnalisé de la classe
\newcommand{\UPSTIclasse}{SNT}

% Type de document
% 0: Custom*				7: Fiche Méthode			14: Document Réponses
% 1: Cours (par défaut)		8: Fiche Synthèse    		15: Programme de colle 
% 2: TD     				9: Formulaire
% 3: TP						10: Memo
% 4: Colle					11: Dossier Technique
% 5: DS						12: Dossier Ressource
% 6: DM						13: Concours Blanc
% * Si on met la valeur 0, il faut décommenter la ligne suivante: 		
%\newcommand{\UPSTItypeDocument}{Custom}
\newcommand{\UPSTIidTypeDocument}{3}

% Titre dans l'en-tête
\newcommand{\UPSTItitreEnTete}{LE LANGAGE HTML}      
\newcommand{\UPSTItitreEnTetePages}{Site web marchand}      
\newcommand{\UPSTIsousTitreEnTete}{Site web marchand}      

% Version du document
\newcommand{\UPSTInumeroVersion}{1.0}
%v1.1 - Ajout de la chaîne fonctionnelle

%----------------------------------------------- 
\UPSTIcompileVars		% "Compile" les variables
%%%%%%%%%%%%%%%%%%%%%%%%%%%%%%%%%%%%%%%%%%%%%%%% 


%%%%%%%%%%%%%%%%%%%%%%%%%%%%%%%%%%%%%%%%%%%%%%%% 
% Début du document
%%%%%%%%%%%%%%%%%%%%%%%%%%%%%%%%%%%%%%%%%%%%%%%% 
\begin{document}
\UPSTIbuildPage

\UPSTIobjectif{Sur cette première activité vous allez modifier le code HTML du site, chaque élève sera responsable d'une page. 
}

\competencetable{Compétences à valider}{
\competence{Etudier et modifier une page HTML simple}
}

\section*{REPARTITION DU TRAVAIL}

\activiteressources
{Organiser au sein du groupe le travail à réaliser}
{
\task Choisir un nom de site pour votre boutique ;
\task Répartir dans le groupe les différentes pages à modifier ;
\task Télécharger le dossier et le dézipper dans votre espace réseau.
}
{documentation}
{}


\section*{MODIFICATION DE LA PARTIE HTML}

\activiteressources
{Modifier les différents textes dans les pages web.}
{
\task Modifier le titre de la page, il doit-être commun à l'ensemble des pages ;
\task Modifier l'ensemble des titres de votre page web, ( balises : h1, h2, h3 ... ) ;
\task Modifier les textes descriptifs ( balises <p> ) ;
\task Ajouter le menu.
}
{documentation}
{NotePad ++}

\section*{MISE EN PLACE DES IMAGES}

\activiteressources
{Ajouter le logo et les images nécessaires.}
{%
\task Ajouter l'image de votre logo dans le dossier image ;
\task Ajouter d'autres images si nécessaires.
}
{documentation}
{NotePad ++}

\section*{TESTER LES DIFFERENTS LIENS}

\activiteressources
{Test des liens hypertext}
{
\task Sur un PC du groupe, copier les différentes pages modifiées ;
\task Ajouter aux dossiers images les différentes images nécessaires ;
\task Tester les différents liens.
}
{documentation}
{NotePad ++}



\end{document}
